% Generated draft: requires mathematical and source review.
\documentclass[11pt]{article}
\usepackage[margin=25mm]{geometry}
\usepackage{fontspec}
\setmainfont{Latin Modern Roman}
\usepackage{amsmath,amssymb,mathtools,unicode-math}
\setmathfont{Latin Modern Math}
\usepackage{enumitem,needspace,xurl,hyperref}
\setlength{\parindent}{0pt}
\setlength{\parskip}{5pt}
\setlength{\emergencystretch}{3em}
\newcommand{\N}{\mathbb{N}}
\newcommand{\Z}{\mathbb{Z}}
\newcommand{\Q}{\mathbb{Q}}
\newcommand{\R}{\mathbb{R}}
\newcommand{\C}{\mathbb{C}}
\newcommand{\F}{\mathbb{F}}
\newcommand{\A}{\mathbb{A}}
\newcommand{\PP}{\mathbb{P}}
\newcommand{\E}{\mathbb{E}}
\newcommand{\eps}{\varepsilon}
\newcommand{\dd}{\,\mathrm{d}}
\begin{document}
\section*{UnsolvedMath ID: 30002548}
\textbf{Source problem code:} catalogue entry
\subsection{Definitions and mathematical statement}
Let \(G\) be a group with identity \(e\). Its integral group algebra \(\Z G\) consists of finite sums \(u=\sum_{g\in G}a_g g\), where \(a_g\in\Z\), with multiplication induced by the group operation. No positivity or symmetry assumption is imposed on the coefficients. For \(n\geq 1\), define
\[
f_u(n)=[e]u^n,\qquad F_u(x)=\sum_{n\geq 1}f_u(n)x^n\in\Z[[x]],
\]
where \([e]\) extracts the coefficient of the identity; this is the notation \([1]\) in the supplied record.

Use the following characteristic-zero conventions. A series \(F(x)\in\Q[[x]]\) is algebraic over \(\Q(x)\) if there exists a nonzero polynomial \(P(x,y)\in\Q[x,y]\) such that \(P(x,F(x))=0\). It is D-finite if there exist \(r\geq 0\) and polynomials \(p_0(x),\ldots,p_r(x)\in\Q[x]\), not all zero, such that
\[
\sum_{j=0}^{r}p_j(x)\frac{\dd^j F}{\dd x^j}(x)=0
\]
as an identity of formal power series.

The first question asks whether, for \(G=\operatorname{GL}(d,\Z)\), the series \(F_u(x)\) is D-finite for every \(u\in\Z G\). Here \(\operatorname{GL}(d,\Z)\) denotes the group of integer \(d\times d\) matrices with determinant \(\pm1\). The supplied record does not explicitly quantify \(d\) or specify a restricted range. Under the usual interpretation that the question concerns all positive integer matrix dimensions, its formulation is
\[
\forall d\in\Z_{\geq1}\;\forall u\in\Z[\operatorname{GL}(d,\Z)],\qquad F_u(x)\text{ is D-finite}.
\]
This interpretation of the dimension range requires source review.

More generally, characterize the groups \(G\) satisfying each of the following two properties separately:
\[
\forall u\in\Z G,\qquad F_u(x)\text{ is algebraic over }\Q(x),
\]
\[
\forall u\in\Z G,\qquad F_u(x)\text{ is D-finite}.
\]
No finite-generation assumption on \(G\) is supplied. The source record does not explicitly name the coefficient field for algebraicity or D-finiteness; the definitions above state the conventions used in this draft.
\subsection{Short English statement}
For each finite integer linear combination of group elements, count the identity coefficient in every positive power and form its generating series. Ask whether every such series is D-finite for integer general linear groups, and characterize the groups for which all such series are algebraic or all are D-finite. The dimension range and coefficient-field conventions are explicitly qualified above.
\subsection{Sources}
\begin{itemize}
\item Local UnsolvedMath registry, supplied complete record with permanent ID \(30002548\), title \emph{Algebraicity and D-Finiteness of Group-Algebra Series}, source problem code \texttt{catalogue entry}, in the Oberwolfach Reports Open Problems collection.
\item Supplied report citation: \emph{Enumerative Combinatorics} (2014), Oberwolfach Reports, pp.~635--720, DOI \texttt{10.4171/owr/2014/12}. \url{https://doi.org/10.4171/owr/2014/12}. The supplied background describes Oberwolfach Reports as documenting mathematical workshops and their open-problem discussions. No individual proposer is identified in the record.
\item Supplied literature-source label: ``Stanley Oberwolfach report.'' \url{https://ems.press/content/serial-article-files/46502?nt=1}. This label is retained without inferring an individual proposer attribution.
\item Supplied literature assessment, checked on \texttt{2026-08-21}: the general D-finiteness question for identity-coefficient series from group algebras of \(\operatorname{GL}(d,\Z)\) is described as unresolved; the assessment also states that undecidability of vanishing does not settle D-finiteness. These are supplied dated curation claims, not independently verified conclusions. The background explicitly requests independent verification of time-sensitive status and cited claims.
\item The registry labels the problem ``open'' and the statement \texttt{corrected\_verified}. Its supplied verification note says that the statement was checked against the official Oberwolfach report and that the clean statement restores the source's hypotheses and notation. These are registry assertions, not verification performed for this draft.
\end{itemize}
\end{document}
